\documentclass[11pt,a4paper]{article}
\usepackage{fontspec,amsmath,amssymb,amsthm,microtype}
\setmainfont{lmroman10-regular.otf}[BoldFont=lmroman10-bold.otf,ItalicFont=lmroman10-italic.otf,BoldItalicFont=lmroman10-bolditalic.otf]
\usepackage{xcolor}
\usepackage[most]{tcolorbox}
\usepackage{mathrsfs}
\providecommand{\lt}{<}
\providecommand{\gt}{>}
\setmonofont{DejaVuSansMono.ttf}[BoldFont=DejaVuSansMono-Bold.ttf,ItalicFont=DejaVuSansMono-Oblique.ttf,BoldItalicFont=DejaVuSansMono-BoldOblique.ttf,Scale=0.86]
\newfontfamily\leanfallback{STIXTwoMath-Regular.otf}[Scale=0.86]
\definecolor{leanrule}{gray}{0.62}
\definecolor{leanback}{gray}{0.962}
\definecolor{leantitle}{gray}{0.22}
\definecolor{leanmuted}{gray}{0.38}
\newtcolorbox{leanbox}[1]{enhanced,breakable,colback=leanback,colframe=leanrule,
  boxrule=0.4pt,arc=0.6mm,left=2.4mm,right=2.4mm,top=1.8mm,bottom=1.8mm,
  before skip=0.8em,after skip=1em,fonttitle=\small\bfseries,coltitle=leantitle,
  colbacktitle=leanback,titlerule=0.3pt,title={#1}}
\newlength{\leancharwidth}
\newenvironment{leancode}{\par\smallskip\begingroup\ttfamily\footnotesize\raggedright
  \settowidth{\leancharwidth}{M}\setlength{\parskip}{0pt}\setlength{\parindent}{0pt}
  \frenchspacing}{\par\endgroup\smallskip}
\newcommand{\leanline}[2]{\hangindent=\dimexpr#1\leancharwidth+4\leancharwidth\relax
  \hangafter=1\noindent\hspace*{#1\leancharwidth}#2\par}
\newcommand{\leansource}[1]{\par{\footnotesize\color{leanmuted}#1}\par}
\newcommand{\leanlabel}[1]{\par\smallskip\noindent{\small\bfseries #1}\par}
\usepackage[margin=25mm]{geometry}
\usepackage[unicode,hidelinks]{hyperref}
\hypersetup{pdfauthor={LangevinLean Project}}
\setlength{\emergencystretch}{3em}
\setlength{\parskip}{0.4em}
\setlength{\parindent}{0pt}
\allowdisplaybreaks[2]
\raggedbottom
\title{\texorpdfstring{Integral formulas for relative entropy and Fisher information}{Integral formulas for relative entropy and Fisher information}}
\author{LangevinLean Project}
\date{}
\begin{document}\maketitle
Draft mathematical exposition; these notes may contain mistakes.


Let \((E,\mathcal E)\) be a measurable space, \(\nu\) a measure on it, and \(r:E\to\mathbb R\) a measurable function with \(r\ge0\); \(r\nu\) is the measure \(A\mapsto\int_Ar\,d\nu\). The relative entropy \(H_\nu(r\nu)\) is defined through the Radon–Nikodym derivative \(d(r\nu)/d\nu\) and its logarithm, the log-likelihood ratio, while the computations in the entropy analysis are carried out with the function \(r\) itself. This section connects the two descriptions. It proves the change-of-measure formula \(\int g\,d(r\nu)=\int rg\,d\nu\); the equivalence \(\log r\in L^1(r\nu)\iff r\log r\in L^1(\nu)\); the identity of the log-likelihood ratio with \(\log r\) almost everywhere; and, as consequences, that \(r\log r\in L^1(\nu)\) implies \(H_\nu(r\nu)<\infty\) and, for probability measures,


\[H_\nu(r\nu)=\int_Er\log r\,d\nu.\]


When \(E\) is a real inner product space, it also proves that square integrability of \(\|\nabla\log r\|\) under \(r\nu\) gives finite Fisher information with \(J_\nu(r)=\int r\|\nabla\log r\|^2\,d\nu\). Only measurability and nonnegativity of \(r\) are needed for the measure-change facts; positivity and smoothness enter only through the Fisher information.




In this section integrals of real or vector-valued functions are Lebesgue (Bochner) integrals, with the convention that the integral of a function that is not integrable is \(0\); every statement below is either an identity valid under this convention or carries an integrability hypothesis. For a finite measure \(\rho\ll\nu\) with log-likelihood ratio \(\ell=\log\frac{d\rho}{d\nu}\in L^1(\rho)\), the relative entropy is \(H_\nu(\rho)=\int\ell\,d\rho+\nu(E)-\rho(E)\), a nonnegative real number; otherwise \(H_\nu(\rho)=\infty\). For probability measures the correction \(\nu(E)-\rho(E)\) vanishes.


\section{\texorpdfstring{Integration against a measure with density}{Integration against a measure with density}}


\textbf{Claim.} Let \(\nu\) be a measure on \(E\), let \(r:E\to\mathbb R\) be measurable with \(r(x)\ge0\) for all \(x\), and let \(g:E\to\mathbb R\) be any function. Then


\[\int_Eg\,d(r\nu)=\int_Er\,g\,d\nu,\]


with the convention that the integral of a non-integrable function is \(0\).




\textbf{Proof.} \emph{Nonnegative measurable functions.} For a measurable set \(A\), \((r\nu)(A)=\int_Ar\,d\nu=\int r\mathbf 1_A\,d\nu\), so the formula \(\int f\,d(r\nu)=\int rf\,d\nu\) holds for indicator functions, by linearity for nonnegative simple functions, and by the monotone convergence theorem (applied on both sides to an increasing sequence of simple functions converging to \(f\)) for every measurable \(f:E\to[0,\infty]\).




\emph{Integrability.} Applying this to \(f=|g|\) for measurable \(g\) gives \(\int|g|\,d(r\nu)=\int r|g|\,d\nu=\int|rg|\,d\nu\), so \(g\in L^1(r\nu)\) if and only if \(rg\in L^1(\nu)\). Measurability up to null sets transfers as well: if \(g\) agrees \(r\nu\)-almost everywhere with a measurable \(\tilde g\), then \(rg=r\tilde g\) \(\nu\)-almost everywhere, because the two can differ only on a \(\nu\)-null subset of \(\{r>0\}\) or on \(\{r=0\}\), where both vanish; conversely, if \(rg=h\) \(\nu\)-almost everywhere with \(h\) measurable, then \(g=\mathbf 1_{\{r>0\}}h/r\) holds \(r\nu\)-almost everywhere, since \((r\nu)(\{r=0\})=0\).




Thus \(g\) is integrable with respect to \(r\nu\) if and only if \(rg\) is integrable with respect to \(\nu\).




\emph{The identity.} If \(g\in L^1(r\nu)\), write \(g=g^+-g^-\); both parts are integrable, the nonnegative case applies to each, and subtracting gives \(\int g\,d(r\nu)=\int rg^+\,d\nu-\int rg^-\,d\nu=\int rg\,d\nu\). If \(g\notin L^1(r\nu)\), then \(rg\notin L^1(\nu)\) and both sides are \(0\) by the convention. \(\square\)




The density is finite at every point, so the product \(rg\) is an ordinary product of real numbers. This formula is used for \(g=\log r\) and \(g=\|\nabla\log r\|^2\) below.

\begin{leanbox}{Lean: \texttt{integral\_\allowbreak{}tiltedMeasure}}
\leansource{\texttt{integral\_\allowbreak{}tiltedMeasure} \quad UnadjustedLangevin/\allowbreak{}LSI/\allowbreak{}KLOfDensity.lean, lines 27–33}
\begin{leancode}
\leanline{0}{theorem\ integral\_tiltedMeasure\ (ν\ :\ Measure\ E)\ (r\ :\ E\ →\ ℝ)}
\leanline{4}{(hmeas\ :\ Measurable\ r)\ (hr\ :\ ∀\ x,\ 0\ ≤\ r\ x)\ (g\ :\ E\ →\ ℝ)\ :}
\leanline{4}{(∫\ x,\ g\ x\ ∂tiltedMeasure\ ν\ r)\ =\ ∫\ x,\ r\ x\ *\ g\ x\ ∂ν}
\end{leancode}
\leanlabel{How the Lean proof follows the argument}
\texttt{integral\_\allowbreak{}tilted\allowbreak{}Measure} has hypotheses \texttt{hmeas : Measurable r} and \texttt{hr : \ensuremath{\forall} x, 0 \ensuremath{\le} r x} and no hypothesis on \texttt{g}; both sides are Bochner integrals \texttt{\ensuremath{\int} x, .\allowbreak{}.\allowbreak{}.\allowbreak{} \ensuremath{\partial}\ensuremath{\mu}}, which Mathlib defines to be \(0\) for non-integrable functions, so the statement is exactly the claim with its convention. After \texttt{unfold tilted\allowbreak{}Measure} the left side is an integral against \texttt{\ensuremath{\nu}.\allowbreak{}with\allowbreak{}Density (fun x \ensuremath{\mapsto} ENNReal.\allowbreak{}of\allowbreak{}Real (r x))}. Mathlib's \texttt{integral\_\allowbreak{}with\allowbreak{}Density\_\allowbreak{}eq\_\allowbreak{}integral\_\allowbreak{}to\allowbreak{}Real\_\allowbreak{}smul} turns it into \(\int f(x)\cdot g(x)\,d\nu\) with \(f\) the real value of the density; its hypotheses are measurability of the density (\texttt{hmeas.\allowbreak{}ennreal\_\allowbreak{}of\allowbreak{}Real}) and finiteness of the density almost everywhere, here at every point (\texttt{ENNReal.\allowbreak{}of\allowbreak{}Real\_\allowbreak{}lt\_\allowbreak{}top}). The final \texttt{simp only [ENNReal.\allowbreak{}to\allowbreak{}Real\_\allowbreak{}of\allowbreak{}Real (hr \_\allowbreak{}), smul\_\allowbreak{}eq\_\allowbreak{}mul]} replaces the real value of \(\max\{r,0\}\) by \(r\) (since \(r\ge0\)) and scalar multiplication by multiplication.
\end{leanbox}

\section{\texorpdfstring{Two equivalent integrability conditions for the logarithm}{Two equivalent integrability conditions for the logarithm}}


\textbf{Claim.} Let \(\nu\) be a measure on \(E\) and \(r:E\to\mathbb R\) measurable with \(r\ge0\). Then


\[\log r\in L^1(r\nu)\quad\Longleftrightarrow\quad r\log r\in L^1(\nu).\]


Here \(\log0\) is given the value \(0\); this value plays no role, since \((r\nu)(\{r=0\})=0\) and \(r\log r=0\) on \(\{r=0\}\).




\textbf{Proof.} This is the integrability part of the previous step with \(g=\log r\), which is measurable: \(\int|\log r|\,d(r\nu)=\int r|\log r|\,d\nu\), so one side is finite exactly when the other is. \(\square\)




The equivalence concerns the absolute integrals, so it controls the positive and the negative parts of the logarithm simultaneously. It converts the integrability of the log-likelihood ratio required by the definition of relative entropy into a condition on the entropy integrand \(r\log r\) with respect to the reference measure.

\begin{leanbox}{Lean: \texttt{integrable\_\allowbreak{}log\_\allowbreak{}tiltedMeasure\_\allowbreak{}iff}}
\leansource{\texttt{integrable\_\allowbreak{}log\_\allowbreak{}tiltedMeasure\_\allowbreak{}iff} \quad UnadjustedLangevin/\allowbreak{}LSI/\allowbreak{}KLOfDensity.lean, lines 37–45}
\begin{leancode}
\leanline{0}{theorem\ integrable\_log\_tiltedMeasure\_iff\ (ν\ :\ Measure\ E)\ (r\ :\ E\ →\ ℝ)}
\leanline{4}{(hmeas\ :\ Measurable\ r)\ (hr\ :\ ∀\ x,\ 0\ ≤\ r\ x)\ :}
\leanline{4}{Integrable\ (fun\ x\ ↦\ Real.log\ (r\ x))\ (tiltedMeasure\ ν\ r)\ ↔}
\leanline{6}{Integrable\ (fun\ x\ ↦\ r\ x\ *\ Real.log\ (r\ x))\ ν}
\end{leancode}
\leanlabel{How the Lean proof follows the argument}
\texttt{integrable\_\allowbreak{}log\_\allowbreak{}tilted\allowbreak{}Measure\_\allowbreak{}iff} applies Mathlib's \texttt{integrable\_\allowbreak{}with\allowbreak{}Density\_\allowbreak{}iff\_\allowbreak{}integrable\_\allowbreak{}smul'} (integrability of \(g\) with respect to \(f\nu\) is equivalent to integrability of \(f\cdot g\) with respect to \(\nu\), for a measurable density \(f\) finite almost everywhere) to \(g=\log r\) and the density \texttt{ENNReal.\allowbreak{}of\allowbreak{}Real (r x)}; \texttt{simpa only [ENNReal.\allowbreak{}to\allowbreak{}Real\_\allowbreak{}of\allowbreak{}Real (hr \_\allowbreak{}), smul\_\allowbreak{}eq\_\allowbreak{}mul]} simplifies the weight to \(r\). Mathlib's \texttt{Real.\allowbreak{}log} satisfies \(\log0=0\), as in the claim.
\end{leanbox}

\section{\texorpdfstring{The log-likelihood ratio of a measure with density}{The log-likelihood ratio of a measure with density}}


\textbf{Claim.} Let \(\nu\) be a \(\sigma\)-finite measure on \(E\) and \(r:E\to\mathbb R\) measurable with \(r\ge0\). Then the log-likelihood ratio of \(r\nu\) with respect to \(\nu\) satisfies


\[\log\frac{d(r\nu)}{d\nu}(x)=\log r(x)\qquad\text{for }(r\nu)\text{-almost every }x.\]


\textbf{Proof.} If \(\nu(A)=0\) then \((r\nu)(A)=\int_Ar\,d\nu=0\), so \(r\nu\ll\nu\). Since \(\nu\) is \(\sigma\)-finite, the Radon–Nikodym derivative of the measure \(r\nu\) with respect to \(\nu\) exists and is \(\nu\)-almost everywhere unique; as \(r\) is a measurable density of \(r\nu\), we get \(\frac{d(r\nu)}{d\nu}=r\) \(\nu\)-almost everywhere. The exceptional set is \(\nu\)-null, hence \((r\nu)\)-null by absolute continuity, so the equality also holds \((r\nu)\)-almost everywhere. Taking logarithms of the two equal nonnegative real numbers gives the claim. \(\square\)




The log-likelihood ratio is defined only up to null sets, through an arbitrary version of the Radon–Nikodym derivative; this claim allows it to be replaced by the explicit function \(\log r\) inside integrals with respect to \(r\nu\).

\begin{leanbox}{Lean: \texttt{llr\_\allowbreak{}tiltedMeasure\_\allowbreak{}ae\_\allowbreak{}eq\_\allowbreak{}log}}
\leansource{\texttt{llr\_\allowbreak{}tiltedMeasure\_\allowbreak{}ae\_\allowbreak{}eq\_\allowbreak{}log} \quad UnadjustedLangevin/\allowbreak{}LSI/\allowbreak{}KLOfDensity.lean, lines 49–55}
\begin{leancode}
\leanline{0}{theorem\ llr\_tiltedMeasure\_ae\_eq\_log\ (ν\ :\ Measure\ E)\ [SigmaFinite\ ν]}
\leanline{4}{(r\ :\ E\ →\ ℝ)\ (hmeas\ :\ Measurable\ r)\ (hr\ :\ ∀\ x,\ 0\ ≤\ r\ x)\ :}
\leanline{4}{llr\ (tiltedMeasure\ ν\ r)\ ν\ =ᵐ[tiltedMeasure\ ν\ r]\ fun\ x\ ↦\ Real.log\ (r\ x)}
\end{leancode}
\leanlabel{How the Lean proof follows the argument}
In \texttt{llr\_\allowbreak{}tilted\allowbreak{}Measure\_\allowbreak{}ae\_\allowbreak{}eq\_\allowbreak{}log}, \texttt{[Sigma\allowbreak{}Finite \ensuremath{\nu}]} is the \(\sigma\)-finiteness assumption, \texttt{llr \ensuremath{\mu} \ensuremath{\nu}} is Mathlib's log-likelihood ratio \(x\mapsto\log\) of the real value of \texttt{\ensuremath{\mu}.\allowbreak{}rn\allowbreak{}Deriv \ensuremath{\nu} x} (the chosen version of \(d\mu/d\nu\), with values in \texttt{\ensuremath{\mathbb R}\ensuremath{\ge}0\ensuremath{\infty}}), and \texttt{f =ᵐ[\ensuremath{\mu}] g} means \(f=g\) \(\mu\)-almost everywhere. \texttt{hac} is the absolute continuity \(r\nu\ll\nu\), from \texttt{with\allowbreak{}Density\_\allowbreak{}absolutely\allowbreak{}Continuous}. \texttt{h\allowbreak{}RN} is \texttt{Measure.\allowbreak{}rn\allowbreak{}Deriv\_\allowbreak{}with\allowbreak{}Density}: the Radon–Nikodym derivative of \(f\nu\) is \(f\), \(\nu\)-almost everywhere. \texttt{hac.\allowbreak{}ae\_\allowbreak{}le h\allowbreak{}RN} transfers this almost-everywhere statement from \(\nu\) to \(r\nu\), and \texttt{filter\_\allowbreak{}upwards} fixes a point \texttt{x} of the resulting full-measure set, where \texttt{simp only [llr, tilted\allowbreak{}Measure, hx, ENNReal.\allowbreak{}to\allowbreak{}Real\_\allowbreak{}of\allowbreak{}Real (hr x)]} computes \(\log r(x)\).
\end{leanbox}

\section{\texorpdfstring{Finite entropy from an integrable entropy integrand}{Finite entropy from an integrable entropy integrand}}


\textbf{Claim.} Let \(\nu\) be \(\sigma\)-finite, let \(r:E\to\mathbb R\) be measurable with \(r\ge0\), and assume \(r\log r\in L^1(\nu)\). Then \(H_\nu(r\nu)<\infty\).




\textbf{Proof.} The divergence \(H_\nu(r\nu)\) is finite as soon as \(r\nu\ll\nu\) and the log-likelihood ratio \(\ell=\log\frac{d(r\nu)}{d\nu}\) is integrable with respect to \(r\nu\). Absolute continuity was shown in the previous step. By the second step, \(r\log r\in L^1(\nu)\) gives \(\log r\in L^1(r\nu)\). By the previous step \(\ell=\log r\) \((r\nu)\)-almost everywhere, and a function that agrees almost everywhere with an integrable function is integrable. Hence \(\ell\in L^1(r\nu)\) and \(H_\nu(r\nu)<\infty\). \(\square\)




No normalization of \(r\nu\) or \(\nu\) is needed for finiteness. This criterion is how finite entropy is established for the explicit densities of the interpolation: one proves integrability of \(r\log r\), a condition on a concrete function, rather than integrability of an abstract log-likelihood ratio.

\begin{leanbox}{Lean: \texttt{klDiv\_\allowbreak{}tiltedMeasure\_\allowbreak{}ne\_\allowbreak{}top}}
\leansource{\texttt{klDiv\_\allowbreak{}tiltedMeasure\_\allowbreak{}ne\_\allowbreak{}top} \quad UnadjustedLangevin/\allowbreak{}LSI/\allowbreak{}KLOfDensity.lean, lines 59–66}
\begin{leancode}
\leanline{0}{theorem\ klDiv\_tiltedMeasure\_ne\_top\ (ν\ :\ Measure\ E)\ [SigmaFinite\ ν]}
\leanline{4}{(r\ :\ E\ →\ ℝ)\ (hmeas\ :\ Measurable\ r)\ (hr\ :\ ∀\ x,\ 0\ ≤\ r\ x)}
\leanline{4}{(hint\ :\ Integrable\ (fun\ x\ ↦\ r\ x\ *\ Real.log\ (r\ x))\ ν)\ :}
\leanline{4}{InformationTheory.klDiv\ (tiltedMeasure\ ν\ r)\ ν\ ≠\ ⊤}
\end{leancode}
\leanlabel{How the Lean proof follows the argument}
\texttt{kl\allowbreak{}Div\_\allowbreak{}tilted\allowbreak{}Measure\_\allowbreak{}ne\_\allowbreak{}top} concludes \texttt{Information\allowbreak{}Theory.\allowbreak{}kl\allowbreak{}Div (tilted\allowbreak{}Measure \ensuremath{\nu} r) \ensuremath{\nu} ≠ \ensuremath{\top}}. It applies Mathlib's criterion \texttt{Information\allowbreak{}Theory.\allowbreak{}kl\allowbreak{}Div\_\allowbreak{}ne\_\allowbreak{}top}, whose two hypotheses are absolute continuity (\texttt{with\allowbreak{}Density\_\allowbreak{}absolutely\allowbreak{}Continuous}) and integrability of \texttt{llr} under the first measure. The latter is \texttt{((integrable\_\allowbreak{}log\_\allowbreak{}tilted\allowbreak{}Measure\_\allowbreak{}iff \ensuremath{\nu} r hmeas hr).\allowbreak{}mpr hint).\allowbreak{}congr (llr\_\allowbreak{}tilted\allowbreak{}Measure\_\allowbreak{}ae\_\allowbreak{}eq\_\allowbreak{}log \ensuremath{\nu} r hmeas hr).\allowbreak{}symm}: integrability of \(\log r\) under \(r\nu\), transported along the almost-everywhere equality \(\log r=\ell\) by \texttt{Integrable.\allowbreak{}congr}.
\end{leanbox}

\section{\texorpdfstring{The entropy as an integral of the density}{The entropy as an integral of the density}}


\textbf{Claim.} Let \(\nu\) be a probability measure, let \(r:E\to\mathbb R\) be measurable with \(r\ge0\), assume that \(r\nu\) is a probability measure and that \(r\log r\in L^1(\nu)\). Then \(H_\nu(r\nu)\) is finite and


\[H_\nu(r\nu)=\int_Er\log r\,d\nu.\]


\textbf{Proof.} Finiteness is the previous step. Let \(\ell\) be the log-likelihood ratio; as in the previous step, \(r\nu\ll\nu\) and \(\ell\in L^1(r\nu)\), so by definition


\[H_\nu(r\nu)=\int\ell\,d(r\nu)+\nu(E)-(r\nu)(E)=\int\ell\,d(r\nu)+1-1=\int\ell\,d(r\nu).\]


Since \(\ell=\log r\) \((r\nu)\)-almost everywhere (third step), \(\int\ell\,d(r\nu)=\int\log r\,d(r\nu)\), and by the change-of-measure formula of the first step with \(g=\log r\), \(\int\log r\,d(r\nu)=\int r\log r\,d\nu\). \(\square\)




Both probability assumptions are used to cancel the correction term \(\nu(E)-(r\nu)(E)\) in the definition of the divergence for finite measures. The resulting signed integral is what is differentiated in time in the entropy calculation; its use requires the integrability hypothesis, which also guarantees that the extended divergence is finite.

\begin{leanbox}{Lean: \texttt{densityEntropy\_\allowbreak{}eq\_\allowbreak{}integral}}
\leansource{\texttt{densityEntropy\_\allowbreak{}eq\_\allowbreak{}integral} \quad UnadjustedLangevin/\allowbreak{}LSI/\allowbreak{}KLOfDensity.lean, lines 70–84}
\begin{leancode}
\leanline{0}{theorem\ densityEntropy\_eq\_integral\ (ν\ :\ Measure\ E)\ [IsProbabilityMeasure\ ν]}
\leanline{4}{(r\ :\ E\ →\ ℝ)\ (hmeas\ :\ Measurable\ r)\ (hr\ :\ ∀\ x,\ 0\ ≤\ r\ x)}
\leanline{4}{(hprob\ :\ IsProbabilityMeasure\ (tiltedMeasure\ ν\ r))}
\leanline{4}{(hint\ :\ Integrable\ (fun\ x\ ↦\ r\ x\ *\ Real.log\ (r\ x))\ ν)\ :}
\leanline{4}{densityEntropy\ ν\ r\ =\ ∫\ x,\ r\ x\ *\ Real.log\ (r\ x)\ ∂ν}
\end{leancode}
\leanlabel{How the Lean proof follows the argument}
In \texttt{density\allowbreak{}Entropy\_\allowbreak{}eq\_\allowbreak{}integral} the conclusion \texttt{density\allowbreak{}Entropy \ensuremath{\nu} r = \ensuremath{\int} x, r x * Real.\allowbreak{}log (r x) \ensuremath{\partial}\ensuremath{\nu}} concerns the real entropy (the real value of the divergence). The proof first registers \texttt{hprob} as an instance (\texttt{let := hprob}) and then collects \texttt{hlog} (integrability of \(\log r\) under \(r\nu\)), \texttt{hllr} (the almost-everywhere identity of the third step), \texttt{hllr\_\allowbreak{}int} (integrability of \texttt{llr}) and \texttt{hac} (absolute continuity). After \texttt{unfold density\allowbreak{}Entropy}, Mathlib's \texttt{Information\allowbreak{}Theory.\allowbreak{}to\allowbreak{}Real\_\allowbreak{}kl\allowbreak{}Div hac hllr\_\allowbreak{}int} rewrites the real value of the divergence as \texttt{\ensuremath{\int} a, llr \ensuremath{\mu} \ensuremath{\nu} a \ensuremath{\partial}\ensuremath{\mu} + \ensuremath{\nu}.\allowbreak{}real univ -\allowbreak{} \ensuremath{\mu}.\allowbreak{}real univ}; \texttt{prob\allowbreak{}Real\_\allowbreak{}univ} evaluates both masses to \(1\) and \texttt{add\_\allowbreak{}sub\_\allowbreak{}cancel\_\allowbreak{}right} removes them. \texttt{integral\_\allowbreak{}congr\_\allowbreak{}ae hllr} replaces \texttt{llr} by \(\log r\), and \texttt{integral\_\allowbreak{}tilted\allowbreak{}Measure} changes the measure.
\end{leanbox}

\section{\texorpdfstring{Finite Fisher information from integrability}{Finite Fisher information from integrability}}


\textbf{Claim.} Let \(E\) be a real inner product space with a \(\sigma\)-algebra, \(\nu\) a measure on \(E\) and \(r:E\to\mathbb R\) a function such that \(\|\nabla\log r\|^2\) is integrable with respect to \(r\nu\). Then \(J_\nu(r)<\infty\).




\textbf{Proof.} By definition \(J_\nu(r)\) is the integral, in \([0,\infty]\), of the nonnegative function \(\|\nabla\log r\|^2\) with respect to \(r\nu\). A real function \(f\ge0\) is integrable exactly when it is almost-everywhere measurable and its integral in \([0,\infty]\) is finite. The hypothesis therefore gives \(J_\nu(r)<\infty\). \(\square\)




This is the finiteness condition under which the log-Sobolev inequality becomes an inequality between real numbers.

\begin{leanbox}{Lean: \texttt{relativeFisher\_\allowbreak{}ne\_\allowbreak{}top}}
\leansource{\texttt{relativeFisher\_\allowbreak{}ne\_\allowbreak{}top} \quad UnadjustedLangevin/\allowbreak{}LSI/\allowbreak{}KLOfDensity.lean, lines 89–95}
\begin{leancode}
\leanline{0}{theorem\ relativeFisher\_ne\_top\ (ν\ :\ Measure\ E)\ (r\ :\ E\ →\ ℝ)}
\leanline{4}{(hint\ :\ Integrable}
\leanline{6}{(fun\ x\ ↦\ ‖gradient\ (fun\ y\ ↦\ Real.log\ (r\ y))\ x‖\ \textasciicircum{}\ 2)\ (tiltedMeasure\ ν\ r))\ :}
\leanline{4}{relativeFisher\ ν\ r\ ≠\ ⊤}
\end{leancode}
\leanlabel{How the Lean proof follows the argument}
\texttt{relative\allowbreak{}Fisher\_\allowbreak{}ne\_\allowbreak{}top} takes the integrability hypothesis \texttt{hint} and concludes \texttt{relative\allowbreak{}Fisher \ensuremath{\nu} r ≠ \ensuremath{\top}}. \texttt{hint.\allowbreak{}has\allowbreak{}Finite\allowbreak{}Integral} extracts the finiteness of the integral of the absolute value, and Mathlib's \texttt{has\allowbreak{}Finite\allowbreak{}Integral\_\allowbreak{}iff\_\allowbreak{}of\allowbreak{}Real} (for an almost-everywhere nonnegative real function \(f\), finiteness of \(\int|f|\) is equivalent to \texttt{\ensuremath{\int}\ensuremath{{}^{-}} a, ENNReal.\allowbreak{}of\allowbreak{}Real (f a) \ensuremath{\partial}\ensuremath{\mu} < \ensuremath{\infty}}) converts it into the strict inequality \(J_\nu(r)<\infty\); nonnegativity is \texttt{sq\_\allowbreak{}nonneg}, and \texttt{ne\_\allowbreak{}of\_\allowbreak{}lt} turns \(<\infty\) into \(\neq\infty\).
\end{leanbox}

\section{\texorpdfstring{The Fisher information as an integral against the reference measure}{The Fisher information as an integral against the reference measure}}


\textbf{Claim.} Let \(E\) be a real inner product space with a \(\sigma\)-algebra, \(\nu\) a measure on \(E\), \(r:E\to\mathbb R\) measurable with \(r\ge0\), and assume that \(\|\nabla\log r\|^2\) is integrable with respect to \(r\nu\). Then the finite number \(J_\nu(r)\) satisfies


\[J_\nu(r)=\int_E\|\nabla\log r\|^2\,d(r\nu)=\int_Er\,\|\nabla\log r\|^2\,d\nu.\]


\textbf{Proof.} For an integrable nonnegative function, the integral in \([0,\infty]\) is finite and equals the Lebesgue integral as a real number; this gives the first equality, with \(J_\nu(r)\) finite by the previous step. The second is the change-of-measure formula of the first step with \(g=\|\nabla\log r\|^2\). \(\square\)




For a strictly positive differentiable \(r\) the integrand is \(r\|\nabla\log r\|^2=\|\nabla r\|^2/r\). This representation, as an integral with respect to the reference measure \(\nu\), is the form in which the Fisher information appears after integration by parts in the entropy calculation.

\begin{leanbox}{Lean: \texttt{relativeFisher\_\allowbreak{}toReal\_\allowbreak{}eq\_\allowbreak{}integral}}
\leansource{\texttt{relativeFisher\_\allowbreak{}toReal\_\allowbreak{}eq\_\allowbreak{}integral} \quad UnadjustedLangevin/\allowbreak{}LSI/\allowbreak{}KLOfDensity.lean, lines 99–111}
\begin{leancode}
\leanline{0}{theorem\ relativeFisher\_toReal\_eq\_integral\ (ν\ :\ Measure\ E)\ (r\ :\ E\ →\ ℝ)}
\leanline{4}{(hmeas\ :\ Measurable\ r)\ (hr\ :\ ∀\ x,\ 0\ ≤\ r\ x)}
\leanline{4}{(hint\ :\ Integrable}
\leanline{6}{(fun\ x\ ↦\ ‖gradient\ (fun\ y\ ↦\ Real.log\ (r\ y))\ x‖\ \textasciicircum{}\ 2)\ (tiltedMeasure\ ν\ r))\ :}
\leanline{4}{(relativeFisher\ ν\ r).toReal\ =}
\leanline{6}{∫\ x,\ r\ x\ *\ ‖gradient\ (fun\ y\ ↦\ Real.log\ (r\ y))\ x‖\ \textasciicircum{}\ 2\ ∂ν}
\end{leancode}
\leanlabel{How the Lean proof follows the argument}
\texttt{relative\allowbreak{}Fisher\_\allowbreak{}to\allowbreak{}Real\_\allowbreak{}eq\_\allowbreak{}integral} proves \texttt{(relative\allowbreak{}Fisher \ensuremath{\nu} r).\allowbreak{}to\allowbreak{}Real = \ensuremath{\int} x, r x * \ensuremath{\Vert}gradient (fun y \ensuremath{\mapsto} Real.\allowbreak{}log (r y)) x\ensuremath{\Vert} \textasciicircum{} 2 \ensuremath{\partial}\ensuremath{\nu}} by a two-line \texttt{calc}. The first line is Mathlib's \texttt{integral\_\allowbreak{}eq\_\allowbreak{}lintegral\_\allowbreak{}of\_\allowbreak{}nonneg\_\allowbreak{}ae} read backwards: for an almost-everywhere nonnegative, almost-everywhere strongly measurable \(f\), \(\int f\,d\mu\) equals the real value of \texttt{\ensuremath{\int}\ensuremath{{}^{-}} a, ENNReal.\allowbreak{}of\allowbreak{}Real (f a) \ensuremath{\partial}\ensuremath{\mu}}; nonnegativity is \texttt{sq\_\allowbreak{}nonneg} and measurability is \texttt{hint.\allowbreak{}aestrongly\allowbreak{}Measurable}. The second line is \texttt{integral\_\allowbreak{}tilted\allowbreak{}Measure}.
\end{leanbox}

\end{document}
